% =============================================================
% LATVIA UNDER SOVIET RULE: A BRIEF ACCOUNT
% =============================================================
\section{Latvia under Soviet rule: a brief account}
\label{sec:latvia}

The companion papers have examined the Lithuanian and Estonian series at length; this section reads the Latvian one only far enough to establish what distinguishes it, before the three are compared. Table~\ref{tab:sectoral} gives the growth of each sector by sub-period and Table~\ref{tab:sectoralidx} the indices themselves at selected years; the annual series is in the appendix.

% tab:sectoral
\begin{table}[htbp]
\centering\scriptsize\setlength{\tabcolsep}{3pt}
\caption{Average annual growth of sectoral value-added volume indices, Latvia (per cent)}
\label{tab:sectoral}
\begin{threeparttable}
\begin{tabular}{lrrrrrrcc}
\toprule
Sector & 1919--29 & 1929--38 & 1938--50 & 1950--68 & 1968--89 & 1989--91 & 1938 level & Peak \\
& & & & & & & regained in & year \\
\midrule
Crops & 11.8 & 5.8 & $-$1.5 & 1.5 & 0.3 & $-$10.6 & 1961 & 1984 \\
Livestock & 3.4 & 3.8 & $-$5.1 & 4.0 & 1.7 & $-$5.5 & 1966 & 1988 \\
Forestry and fishing & 10.5 & 0.3 & $-$1.4 & 4.9 & 0.2 & $-$9.8 & 1958 & 1971 \\
\quad Primary sector & 5.7 & 4.3 & $-$3.3 & 2.8 & 1.0 & $-$7.6 & 1967 & 1989 \\
\midrule
Food & --- & --- & $-$1.9 & 7.3 & 2.1 & $-$4.3 & 1953 & 1980 \\
Textiles & --- & --- & 1.7 & 7.2 & 0.6 & $-$12.8 & 1949 & 1987 \\
Other light industries & --- & --- & 4.4 & 6.3 & $-$0.8 & $-$3.7 & 1948 & 1971 \\
Heavy industries & --- & --- & 5.4 & 10.3 & 1.2 & $-$9.1 & 1948 & 1986 \\
\quad Industry & 20.4 & 6.8 & 3.9 & 8.1 & 0.8 & $-$8.0 & 1948 & 1986 \\
\midrule
Construction & 7.8 & 2.1 & 7.4 & 8.4 & $-$0.5 & $-$9.0 & 1947 & 1975 \\
Transport and communications & 21.2 & $-$0.5 & $-$1.4 & 12.2 & 3.2 & $-$37.6 & 1952 & 1988 \\
Trade, accommodation and catering & 14.9 & 3.1 & 1.4 & 10.2 & 4.5 & 6.5 & 1949 & 1990 \\
Finance and insurance & 37.3 & $-$1.9 & $-$0.9 & 7.1 & 5.5 & 9.5 & 1951 & 1990 \\
Housing & 2.2 & 4.2 & $-$5.7 & 6.0 & 2.0 & $-$3.5 & 1959 & 1987 \\
Public services & $-$1.2 & 0.8 & 8.0 & 1.0 & 0.5 & $-$6.2 & 1945 & 1986 \\
\midrule
\textbf{GDP} & 7.4 & 3.6 & 1.2 & 5.2 & 1.4 & $-$8.5 & 1948 & 1986 \\
\bottomrule
\end{tabular}
\begin{tablenotes}[para,flushleft]
\footnotesize
\textit{Notes:} compound annual growth between end-point years of constant-price value-added indices (1938 = 100). Industrial branch indices are available from 1938 only. `1938 level regained': first post-war year (1945 or later) in which the index is at or above its 1938 value; `Peak year': year of the highest index value up to 1990. Public services comprise administration, defence, health and education. Growth in 1938--50 is affected by the wartime trough of 1944, for which Latvian sectoral movements are partly extrapolated from Lithuanian ratios (Section~\ref{sec:method}). The 1989--91 column ends in 1991, the splicing year beyond which the final series follows adjusted Maddison Project data. \textit{Sources:} see online appendix.
\end{tablenotes}
\end{threeparttable}
\end{table}


% tab:sectoralidx
\begin{table}[htbp]
\centering\footnotesize
\caption{Volume indices of sectoral value-added in Latvia, selected years, 1919--95 (1938 = 100)}
\label{tab:sectoralidx}
\begin{threeparttable}
\begin{tabular}{lrrrrrr}
\toprule
Year & GDP & Primary & Industry & Construction & Market & Public \\
& & sector & & & services & services \\
\midrule
1919 & 36 & 39 & 9 & 39 & 27 & 105 \\
1929 & 73 & 68 & 55 & 83 & 81 & 93 \\
1938 & 100 & 100 & 100 & 100 & 100 & 100 \\
\midrule
1945 & 68 & 47 & 27 & 38 & 55 & 200 \\
1950 & 115 & 67 & 159 & 235 & 74 & 253 \\
1955 & 140 & 60 & 293 & 420 & 126 & 219 \\
1960 & 190 & 92 & 433 & 625 & 185 & 207 \\
1968 & 287 & 109 & 645 & 995 & 366 & 305 \\
1975 & 356 & 111 & 814 & 1,183 & 499 & 387 \\
1980 & 365 & 109 & 786 & 976 & 634 & 378 \\
1985 & 392 & 127 & 829 & 962 & 677 & 413 \\
1989 & 383 & 135 & 763 & 900 & 745 & 336 \\
\midrule
1991 & 321 & 115 & 645 & 746 & 592 & 296 \\
1995 & 233 & 63 & 178 & 242 & 745 & 270 \\
\bottomrule
\end{tabular}
\begin{tablenotes}[para,flushleft]
\footnotesize
\textit{Notes:} market services: transport and communications, trade, finance and insurance and housing; public services: administration, defence, health and education. The annual indices are in Tables~\ref{tab:appsec1} and~\ref{tab:appsec2}. \textit{Sources:} see online appendix.
\end{tablenotes}
\end{threeparttable}
\end{table}


Sovietisation began at once. Within a day of the proclamation of Soviet power in July 1940 land, banks and the larger industrial, trade and transport enterprises were declared state property; by the end of the year 1,477 industrial enterprises producing 84 per cent of output had been nationalised, the lat was exchanged for the rouble at par although free-market prices in Latvia were a fraction of Soviet ones, and in March 1941 the lat was withdrawn without warning and savings above 1,000 lats annulled \parencites[183--8]{krumins_2017_tautsaimnieciba_1940}[32--3]{years_of_dependence}[290, 294]{terleckaspinigai}.\footnote{Misiunas and Taagepera date the abolition of the Baltic currencies to late 1940 \parencite[33]{years_of_dependence}; the lat was declared invalid without warning on 25 March 1941 \parencite[188]{krumins_2017_tautsaimnieciba_1940}.} The military aviation plant that Moscow secretly ordered built in Riga in October 1940, with machine tools requisitioned from VEF and the other large works, foreshadowed the role the republic's industry would play \parencite[194--5]{krumins_2017_tautsaimnieciba_1940}. The deportation of some 15,000--15,600 people on 14 June 1941, the German occupation, the Holocaust, which killed some 66,000 Latvian Jews, the mobilisations of 1943--4 and the flight of 1944 reduced the resident population by \lvPOPchangeWARtwo{} per cent between 1939 and 1945; King's balance of the decade 1939--49 puts the loss at 570,000 people, a quarter of the pre-war population, only a tenth of it from military operations \parencites[53--4]{labsvirs_1959_case}[81--94]{king_1965_economic}[42, 64]{years_of_dependence}. Income per head fell on the present estimates by 28 per cent between 1938 and the trough of 1944, and output by two-fifths, but the figures for 1941--5 are the least secure in the series.

Recovery was fast and, in its sources, revealing. Industry regained its 1938 level in \sregain{Industry} and GDP in \sregain{GDP};\footnote{Misiunas and Taagepera likewise date the recovery of pre-war industrial output in Latvia to 1949 \parencite[113]{years_of_dependence}; Mertelsmann, who dates Estonia's recovery to the second half of the 1950s, holds that Latvia followed the same pattern \parencite[167--9]{mertelsmann_2003_there}.} with a smaller population, income per head regained its 1938 level in \lvREGAINthirtyeight{} and had exceeded it by a fifth by 1950.\footnote{Consumption lagged output: urban Latvian families in 1961 still consumed 30--50 per cent less meat, eggs and dairy produce than before the war, and the recovery of pre-war living standards is commonly dated to the 1960s \parencites[185--6]{years_of_dependence}[38]{puur_2019_turbulent}.} The fourth five-year plan built on the machine-building and electrotechnical base of the interwar republic, which had made nearly a third of the Soviet Union's radios in 1939, and on labour imported from the Russian interior: net immigration of 272,000 in 1946--50 accounted for nine-tenths of population growth, and industrial employment doubled between 1945 and 1950 \parencites[46--7]{eglite_2015_migracija}[233--4]{krumins_2017_sovetizacijas}[244]{vaskela_2014_lietuvos}. Moscow's investment, celebrated in Soviet sources as two billion roubles over the five post-war years, came from Latvia's own revenues: the archival budget accounts show that only about a quarter of the revenue raised in the republic went to its own budget, that military and repressive ministries took nearly half of all spending in Latvia in the late 1940s, and that over 1946--90 as a whole the Union retained 18 per cent of what it collected there, and a further 13.3 billion roubles, 15.7 per cent of all revenue raised there, was spent from the republican budget on Soviet military needs \parencites[251--7]{krumins_2017_sovetizacijas}[9]{pazdere_2015_visparejs}{gatis_compensations}.\footnote{The budget accounts omit implicit transfers through prices; Aage suggests that with cheap energy counted the Baltic republics `probably received positive net transfers from Russia', an answer that, as Viksnins had already noted, depends on the price assigned to energy, and that Samonis had rejected for Lithuania \parencites[102]{aage_1997_public}[252]{viksnins_1986_latvian}[6--7]{samonis_1995_transforming}.}

Collectivisation came in 1949 and, as everywhere in the Baltic, by deportation. On 25 March and the days following, more than 40,000 people, most of them farming families, were sent east;\footnote{Feest's count is 41,445 \parencite[99--100]{feest2014}; for the date see \parencites[172]{strods_2006_nonviolent}[75]{pocius_2019_lietuvos}; the older Western literature put the number at 50,000 or more \parencites[90]{labsvirs_1959_case}[99]{years_of_dependence}.} the kolkhoz share of agricultural land rose from a tenth at the beginning of 1949 to 83 per cent a year later \parencite[94]{labsvirs_1959_case}, the share of farms collectivised to 85 per cent on the count used in Table~\ref{tab:collectivisation}, and to 93 per cent by the end of December 1949 on Krūmiņš's,\footnote{Misiunas and Taagepera also give 93 per cent at the end of 1949 \parencite[102]{years_of_dependence}; Terleckas gives 90 per cent on 1 January 1950 \parencite[225]{terleckas_2026_siferio}.} and the cooperative dairies and credit societies of the interwar republic, already dismantled, were replaced by procurement prices that fell for meat from 16,640 to 680 roubles a tonne \parencites[245--50]{krumins_2017_sovetizacijas}[48]{aizsilnieks_1974_sovietization}. Agricultural output, on the present estimates 59 per cent of its 1938 level in 1947, fluctuated between 60 and 67 per cent of it through 1955; grain yields in the early 1950s were little more than half those of the late 1930s, and milk output less than half; the sector regained its pre-war level only in \sregain{Primary}, crops in \sregain{Crops} and livestock in \sregain{Livestock}, and the kolkhoz households' private plots, a tenth of the sown area, still produced more than half of the republic's milk and meat in 1958, while the kolkhoz markets accounted for more than 60 per cent of urban retail sales of meat and milk products \parencites[219--23, 269--71]{labsvirs_1959_case}[282]{krumins_2017_tautsaimnieciba_1950}.

The two decades to 1968 were the fastest the Latvian economy had yet recorded. Industrial output grew by \sgr{Industry}{1950}{1968} per cent a year, heavy industry by \sgr{Heavy}{1948}{1969} per cent a year over 1948--69, construction by \sgr{Construction}{1950}{1968} per cent and transport by \sgr{Transport}{1950}{1968}; total GDP grew by \tgr{1950}{1968} per cent a year and income per head by \gr{1950}{1968}. The growth was extensive in the Soviet sense, carried by machine-building for the Union market, by the electric trains, radios and telephone equipment of Riga and by a second wave of chemical and fibre plants in the 1960s, and by labour: a net 292,000 immigrants arrived between 1960 and 1985, and Latvians fell from 62 per cent of the population in 1959 to 52 per cent in 1989 \parencites[303]{krumins_2017_tautsaimnieciba_1960}[50]{eglite_2015_migracija}[43--4]{puur_2019_turbulent}. The one attempt to steer the republic's own development came with the sovnarkhoz reform of 1957, which placed most of Latvian industry under a republican council and raised the share of revenue retained in the republic from a quarter to a half; the `national communists' who used it to curb heavy industry and immigration were purged in 1959,\footnote{Berklāvs was dismissed in July 1959 after the visit of a Moscow commission, Lācis and Kalnbērziņš followed by the end of the year, and the purge ran for two and a half years; Pelše's leadership then lifted the limit on immigration \parencites[249--50]{smith_2003_republican}[172]{strods_2006_nonviolent}[487--8]{grybkauskas_2022_economic}.} after which, in the judgement of the literature, the republic's lobbying power in Moscow was weaker than Lithuania's or Estonia's and large all-Union industry dominated \parencites[259--65]{krumins_2017_tautsaimnieciba_1950}[203--5]{king_1965_economic}[142--3, 187, 228]{years_of_dependence}. Two irregularities in the series fall in these years: a contraction of 3.6 per cent in income per head in 1962, and the turning point of 1968, after which growth never recovered. The first was the Latvian instance of the agricultural crisis that struck the whole Union in 1962--3, when a poor harvest, the restrictions imposed on private plots and livestock since 1958 and the exhaustion of the virgin-lands and maize campaigns forced the leadership to raise the retail prices of meat and butter at the end of May 1962, the only retail price rise between 1961 and 1985, to ration bread in many cities and to buy grain abroad in 1963 \parencites[17--19]{belkharoev2024}[111--14]{andreenkov2021}[299]{krumins_2017_tautsaimnieciba_1960}. On the present estimates the Latvian crop index fell by nearly half in 1962 and did not return to its 1960 level until 1964; Estonia's fell by nearly a fifth, and Estonia's grain deliveries to the all-Union fund fell from 30,000 to 13,000 tons \parencite[55]{purre1964}. Industrial growth paused in the same year.

The years from 1968 to 1989 brought growth slower than any recorded for the communist bloc, on comparator series that, for East-Central Europe, almost certainly overstate it. Income per head grew by \gr{1968}{1989} per cent a year, and after 1973 by \gr{1973}{1989}; industry grew by \sgr{Industry}{1968}{1989} per cent a year, and in 1970--85 its growth was the lowest of any Soviet republic \parencite[304]{krumins_2017_tautsaimnieciba_1960}. Construction, which had grown by \sgr{Construction}{1950}{1968} per cent a year to 1968, peaked in \speak{Construction} and contracted thereafter; on the present estimates the harvests fell into a depression between 1974 and 1982 that has no counterpart in Estonia or Lithuania; and by the mid-1980s the republic consumed nearly twice the electricity it generated and imported feed grain equal to its own production \parencite[302, 307]{krumins_2017_tautsaimnieciba_1960}. Enterprises of the all-Union ministries produced 34 per cent of industrial output in 1970 and 37 per cent in 1984, and military production, concentrated in VEF, Alfa and Kommutators, employed more than 15 per cent of the labour force in 1985, three times the Estonian share \parencites[34]{gulans_2015_rupnieciba}[103--4]{wbtransitionlatv}[59--61]{upmalis_2015_latvijas}. Latvian income per head, \pctussr{1968} per cent of the Soviet average in 1968 on the present estimates, stood at \pctussr{1989} per cent of it by 1989, and the shops of Latvia were worse supplied from the late 1970s than those of Estonia or Lithuania \parencites[304--5]{krumins_2017_tautsaimnieciba_1960}[318]{krumins_2017_gorbacova}[222]{years_of_dependence}[141]{demuth_2000_politics}. The decline began before the Union's: income per head peaked in \lvSOVPEAKyear{} and fell by \gra{1986}{1989} per cent a year to 1989, physical output in 1990 was below that of 1985 in most branches \parencite[348--50]{krumins_2017_gorbacova}, and on the present estimates Estonia overtook Latvia durably in 1989, as it had done only briefly before, in the aftermath of each war and in 1962--4.\footnote{Official series differ on the last point, still placing Latvian income per head above Estonian in 1988--9 \parencites[35]{vanselm_1997_economics}[156]{norkus_2023_baltic}.}

The transition then removed the Soviet market on which the whole structure rested. Income per head fell by \lvTROUGHpct{} per cent between 1989 and \lvTROUGHyear{}, industry by about two-thirds in three years, and machine-building, which in 1990 employed 132,000 people and exported 92 per cent of its output, only 3 per cent of it to the European Union, saw its two largest branches fall to roughly a quarter of their 1990 output by 1995; the bus and railway-wagon works had closed by 1998 and VEF had been broken up \parencites[925--9]{silins_2017_metalapstrade}[130]{soms_2015_latvijas}.\footnote{Official series record a deeper fall: the EBRD puts Latvian GDP in 1994 at 52 per cent of its 1989 level \parencite[13]{kukk_2014_road}.} Prices were freed in 1991--2, the lats introduced in 1993 and pegged in 1994, and privatisation, slower than in Estonia and Lithuania, had sold about half of former state assets by mid-1998; the banking crisis of 1995, in which Banka Baltija and twelve other banks failed, cost a further 2--4 per cent of GDP \parencites[811--15]{krastins_2017_kreditiestades}[142]{bitans_2012_development}[205--9]{jeffries_2004_countries}. Recovery began in 1994, and the 1989 level of income per head was regained in \lvREGAINeightynine{}, after Estonia (\eeREGAINeightynine{}) and Lithuania (\ltREGAINeightynine{}).

Three things distinguish the Latvian record. The republic entered the Union with the strongest industrial base of the three, specialised most heavily in all-Union machine-building and military production, and surrendered the largest share of its revenue; it lost its only instrument of economic autonomy in 1959; and it stagnated first, declined earliest, and fell furthest in the transition. The comparison that follows asks how much of this was Latvian and how much was Soviet.
